\chapter{Complexity — Counter Prog Run}
%%% generated by the blueprint pipeline from TCSlib.Complexity.TuringMachine.CounterProgRun
%%%%%%%%%%%%%%%%%%%%%%%%%%%%%%%%%%%%%%%%%%%%%%%%%%%%%%%%%%%%%%%%%%%%%%%%%%%%%%%
\section{Overview}

This module provides Hoare-style reasoning about runs of unary counter programs
(goto programs over natural-number registers that read a binary input once and
print bits or register values in unary), the programming model for the
polynomial-time emitters of Arora--Barak, \S6.2, Remark~6.7.  It introduces the
run relation ``from this label and these registers the program reaches that label
and those registers, printing $e$, within $b$ steps'', its composition, one rule
per instruction, the count-down loop, straight-line templates of print
micro-operations, and linear expressions in the registers printed in unary; it
also shows that a halting run of $t$ abstract steps is a computation of the
compiled multi-tape Turing machine within $t(2t+3)$ steps, and bounds the growth
of the input position and of the output along a run.

%%%%%%%%%%%%%%%%%%%%%%%%%%%%%%%%%%%%%%%%%%%%%%%%%%%%%%%%%%%%%%%%%%%%%%%%%%%%%%%
\section{Declarations}

\begin{definition}[Run relation of a counter program]
\lean{Complexity.CounterProg.Goes}
\label{Complexity.CounterProg.Goes}
\leanok
\uses{Complexity.CounterProg.run}
Let $P$ be a counter program over $R$ natural-number registers with label set
$\Lambda$, run on a binary input $x$.  Given a label $\ell \in \Lambda$, a target
$\ell' \in \Lambda \cup \{\bot\}$ (where $\bot$ means halted), register vectors
$\rho, \rho' : \{0,\dots,R-1\} \to \mathbb{N}$, input positions $p, p' \in
\mathbb{N}$, a bit string $e$ and a bound $b \in \mathbb{N}$, this is the
predicate that for every output $o$ printed so far there is some $t \le b$ such
that $t$ steps of $P$ on $x$ lead from the state with label $\ell$, registers
$\rho$, $p$ input symbols read and output $o$ to the state with label $\ell'$,
registers $\rho'$, $p'$ input symbols read and output $o\,e$.
\end{definition}

\begin{theorem}[Composition of counter-program runs]
\lean{Complexity.CounterProg.Goes.trans}
\label{Complexity.CounterProg.Goes.trans}
\leanok
\uses{Complexity.CounterProg.Goes, Complexity.CounterProg.run_add}
\difficulty{3}
Let $P$ be a counter program on a binary input $x$, let $\ell, \ell'$ be labels,
$\ell''$ a label or the halted state, $\rho, \rho', \rho''$ register vectors and
$p, p', p''$ input positions.  Suppose that from label $\ell$ with registers
$\rho$ and $p$ symbols read, $P$ reaches label $\ell'$ with registers $\rho'$ and
$p'$ symbols read within $b$ steps while appending $e$ to whatever output was
printed before, and that from $(\ell', \rho', p')$ it reaches $(\ell'', \rho'',
p'')$ within $b'$ steps while appending $e'$.  Then from $(\ell, \rho, p)$ it
reaches $(\ell'', \rho'', p'')$ within $b + b'$ steps while appending the
concatenation $e\,e'$.
\end{theorem}

\begin{theorem}[Weakening the step bound of a run]
\lean{Complexity.CounterProg.Goes.mono}
\label{Complexity.CounterProg.Goes.mono}
\leanok
\uses{Complexity.CounterProg.Goes}
\difficulty{2}
Let $P$ be a counter program on a binary input $x$ and let $b \le b'$.  If from
label $\ell$ with registers $\rho$ and $p$ input symbols read, $P$ reaches
$\ell'$ (a label or the halted state) with registers $\rho'$ and $p'$ symbols
read within $b$ steps while appending $e$ to its output, then it does so within
$b'$ steps as well.
\end{theorem}

\begin{theorem}[Rewriting the endpoints of a run]
\lean{Complexity.CounterProg.Goes.congr}
\label{Complexity.CounterProg.Goes.congr}
\leanok
\uses{Complexity.CounterProg.Goes, Complexity.CounterProg.Goes.mono}
\difficulty{2}
Let $P$ be a counter program on a binary input $x$.  Suppose that from label
$\ell$ with registers $\rho$ and $p$ input symbols read, $P$ reaches $\ell'$ (a
label or the halted state) with registers $\rho'$ and $p'$ symbols read within
$b$ steps while appending $e$ to its output, and that $\rho = \sigma$,
$\rho' = \sigma'$, $p = q$, $p' = q'$, $e = e'$ and $b \le b'$.  Then from label
$\ell$ with registers $\sigma$ and $q$ symbols read, $P$ reaches $\ell'$ with
registers $\sigma'$ and $q'$ symbols read within $b'$ steps while appending
$e'$.
\end{theorem}

\begin{theorem}[A single step is a run of length one]
\lean{Complexity.CounterProg.goes_of_step}
\label{Complexity.CounterProg.goes_of_step}
\leanok
\uses{Complexity.CounterProg.Goes, Complexity.CounterProg.step}
\difficulty{1}
Let $P$ be a counter program on a binary input $x$, let $\ell$ be a label,
$\ell'$ a label or the halted state, $\rho, \rho'$ register vectors, $p, p'$
input positions and $e$ a bit string.  If for every output $o$ one step of $P$
on $x$ maps the state with label $\ell$, registers $\rho$, $p$ symbols read and
output $o$ to the state with label $\ell'$, registers $\rho'$, $p'$ symbols read
and output $o\,e$, then from $(\ell, \rho, p)$ the program reaches
$(\ell', \rho', p')$ within one step while appending $e$ to its output.
\end{theorem}

\begin{theorem}[Run rule for the halt instruction]
\lean{Complexity.CounterProg.goes_halt}
\label{Complexity.CounterProg.goes_halt}
\leanok
\uses{Complexity.CounterProg.Goes, Complexity.CounterProg.goes_of_step, Complexity.CounterProg.step}
\difficulty{1}
Let $P$ be a counter program on a binary input $x$ and $\ell$ a label whose
instruction is $\mathsf{halt}$.  Then for all registers $\rho$ and input
positions $p$, from label $\ell$ with registers $\rho$ and $p$ symbols read the
program reaches the halted state, with registers still $\rho$ and still $p$
symbols read, within one step and without printing anything.
\end{theorem}

\begin{theorem}[Run rule for the goto instruction]
\lean{Complexity.CounterProg.goes_goto}
\label{Complexity.CounterProg.goes_goto}
\leanok
\uses{Complexity.CounterProg.Goes, Complexity.CounterProg.goes_of_step, Complexity.CounterProg.step}
\difficulty{1}
Let $P$ be a counter program on a binary input $x$ and $\ell$ a label whose
instruction is $\mathsf{goto}\ \ell'$.  Then for all registers $\rho$ and input
positions $p$, from label $\ell$ with registers $\rho$ and $p$ symbols read the
program reaches label $\ell'$ with registers $\rho$ and $p$ symbols read within
one step, without printing anything.
\end{theorem}

\begin{theorem}[Run rule for printing a bit]
\lean{Complexity.CounterProg.goes_out}
\label{Complexity.CounterProg.goes_out}
\leanok
\uses{Complexity.CounterProg.Goes, Complexity.CounterProg.goes_of_step, Complexity.CounterProg.step}
\difficulty{1}
Let $P$ be a counter program on a binary input $x$ and $\ell$ a label whose
instruction is $\mathsf{out}(b, \ell')$, ``print the bit $b$, then go to
$\ell'$''.  Then for all registers $\rho$ and input positions $p$, from label
$\ell$ with registers $\rho$ and $p$ symbols read the program reaches label
$\ell'$ with registers $\rho$ and $p$ symbols read within one step, appending the
single bit $b$ to its output.
\end{theorem}

\begin{theorem}[Run rule for incrementing a register]
\lean{Complexity.CounterProg.goes_inc}
\label{Complexity.CounterProg.goes_inc}
\leanok
\uses{Complexity.CounterProg.Goes, Complexity.CounterProg.goes_of_step, Complexity.CounterProg.step}
\difficulty{1}
Let $P$ be a counter program on a binary input $x$, $r$ a register and $\ell$ a
label whose instruction is $\mathsf{inc}(r, \ell')$.  Then for all registers
$\rho$ and input positions $p$, from label $\ell$ with registers $\rho$ and $p$
symbols read the program reaches label $\ell'$ within one step, printing
nothing, with $p$ symbols read and registers $\rho$ except that register $r$ now
holds $\rho(r) + 1$.
\end{theorem}

\begin{theorem}[Run rule for decrementing a register]
\lean{Complexity.CounterProg.goes_dec}
\label{Complexity.CounterProg.goes_dec}
\leanok
\uses{Complexity.CounterProg.Goes, Complexity.CounterProg.goes_of_step, Complexity.CounterProg.step}
\difficulty{1}
Let $P$ be a counter program on a binary input $x$, $r$ a register and $\ell$ a
label whose instruction is $\mathsf{dec}(r, \ell')$.  Then for all registers
$\rho$ and input positions $p$, from label $\ell$ with registers $\rho$ and $p$
symbols read the program reaches label $\ell'$ within one step, printing
nothing, with $p$ symbols read and registers $\rho$ except that register $r$ now
holds $\max(\rho(r) - 1, 0)$.
\end{theorem}

\begin{theorem}[Zero test on a zero register]
\lean{Complexity.CounterProg.goes_jz_zero}
\label{Complexity.CounterProg.goes_jz_zero}
\leanok
\uses{Complexity.CounterProg.Goes, Complexity.CounterProg.goes_of_step, Complexity.CounterProg.step}
\difficulty{1}
Let $P$ be a counter program on a binary input $x$, $r$ a register and $\ell$ a
label whose instruction is $\mathsf{jz}(r, \ell_0, \ell_1)$, ``go to $\ell_0$ if
register $r$ is zero, else to $\ell_1$''.  If the registers $\rho$ satisfy
$\rho(r) = 0$, then for every input position $p$, from label $\ell$ with
registers $\rho$ and $p$ symbols read the program reaches label $\ell_0$ with
registers $\rho$ and $p$ symbols read within one step, printing nothing.
\end{theorem}

\begin{theorem}[Zero test on a nonzero register]
\lean{Complexity.CounterProg.goes_jz_pos}
\label{Complexity.CounterProg.goes_jz_pos}
\leanok
\uses{Complexity.CounterProg.Goes, Complexity.CounterProg.goes_of_step, Complexity.CounterProg.step}
\difficulty{1}
Let $P$ be a counter program on a binary input $x$, $r$ a register and $\ell$ a
label whose instruction is $\mathsf{jz}(r, \ell_0, \ell_1)$, ``go to $\ell_0$ if
register $r$ is zero, else to $\ell_1$''.  If the registers $\rho$ satisfy
$\rho(r) \ne 0$, then for every input position $p$, from label $\ell$ with
registers $\rho$ and $p$ symbols read the program reaches label $\ell_1$ with
registers $\rho$ and $p$ symbols read within one step, printing nothing.
\end{theorem}

\begin{theorem}[Run rule for printing a register in unary]
\lean{Complexity.CounterProg.goes_pr}
\label{Complexity.CounterProg.goes_pr}
\leanok
\uses{Complexity.CounterProg.Goes, Complexity.CounterProg.goes_of_step, Complexity.CounterProg.step}
\difficulty{1}
Let $P$ be a counter program on a binary input $x$, $r$ a register and $\ell$ a
label whose instruction is $\mathsf{pr}(r, \ell')$, ``print the value of
register $r$ in unary, then go to $\ell'$''.  Then for all registers $\rho$ and
input positions $p$, from label $\ell$ with registers $\rho$ and $p$ symbols read
the program reaches label $\ell'$ with registers $\rho$ and $p$ symbols read
within one step, appending the string $1^{\rho(r)}$ of $\rho(r)$ ones to its
output.
\end{theorem}

\begin{theorem}[Reading at the end of the input]
\lean{Complexity.CounterProg.goes_rd_end}
\label{Complexity.CounterProg.goes_rd_end}
\leanok
\uses{Complexity.CounterProg.Goes, Complexity.CounterProg.goes_of_step, Complexity.CounterProg.step}
\difficulty{1}
Let $P$ be a counter program on a binary input $x$ and $\ell$ a label whose
instruction is $\mathsf{rd}(\ell_e, \ell_0, \ell_1)$, ``read the next input
symbol and go to $\ell_e$, $\ell_0$ or $\ell_1$ according as the input has ended,
the symbol is $0$, or it is $1$''.  If $x$ has no symbol at (zero-based) position
$p$, i.e.\ $p \ge |x|$, then for all registers $\rho$, from label $\ell$ with
registers $\rho$ and $p$ symbols read the program reaches label $\ell_e$ with
registers $\rho$ and still $p$ symbols read within one step, printing nothing.
\end{theorem}

\begin{theorem}[Reading a zero from the input]
\lean{Complexity.CounterProg.goes_rd_false}
\label{Complexity.CounterProg.goes_rd_false}
\leanok
\uses{Complexity.CounterProg.Goes, Complexity.CounterProg.goes_of_step, Complexity.CounterProg.step}
\difficulty{1}
Let $P$ be a counter program on a binary input $x$ and $\ell$ a label whose
instruction is $\mathsf{rd}(\ell_e, \ell_0, \ell_1)$, ``read the next input
symbol and go to $\ell_e$, $\ell_0$ or $\ell_1$ according as the input has ended,
the symbol is $0$, or it is $1$''.  If the symbol of $x$ at (zero-based) position
$p$ is $0$, then for all registers $\rho$, from label $\ell$ with registers
$\rho$ and $p$ symbols read the program reaches label $\ell_0$ with registers
$\rho$ and $p + 1$ symbols read within one step, printing nothing.
\end{theorem}

\begin{theorem}[Reading a one from the input]
\lean{Complexity.CounterProg.goes_rd_true}
\label{Complexity.CounterProg.goes_rd_true}
\leanok
\uses{Complexity.CounterProg.Goes, Complexity.CounterProg.goes_of_step, Complexity.CounterProg.step}
\difficulty{1}
Let $P$ be a counter program on a binary input $x$ and $\ell$ a label whose
instruction is $\mathsf{rd}(\ell_e, \ell_0, \ell_1)$, ``read the next input
symbol and go to $\ell_e$, $\ell_0$ or $\ell_1$ according as the input has ended,
the symbol is $0$, or it is $1$''.  If the symbol of $x$ at (zero-based) position
$p$ is $1$, then for all registers $\rho$, from label $\ell$ with registers
$\rho$ and $p$ symbols read the program reaches label $\ell_1$ with registers
$\rho$ and $p + 1$ symbols read within one step, printing nothing.
\end{theorem}

\begin{theorem}[The count-down loop]
\lean{Complexity.CounterProg.goes_loop}
\label{Complexity.CounterProg.goes_loop}
\leanok
\uses{Complexity.CounterProg.Goes, Complexity.CounterProg.Goes.congr, Complexity.CounterProg.Goes.trans, Complexity.CounterProg.goes_dec, Complexity.CounterProg.goes_jz_pos, Complexity.CounterProg.goes_jz_zero}
\difficulty{6}
Let $P$ be a counter program on a binary input $x$, $r$ a register, and
$h, d, c, z$ labels such that the instruction at $h$ is
$\mathsf{jz}(r, z, d)$ (go to $z$ if register $r$ is zero, else to $d$) and the
instruction at $d$ is $\mathsf{dec}(r, c)$.  Let $K, b \in \mathbb{N}$, and for
each $j \in \mathbb{N}$ let $\sigma_j$ be a register vector, $p_j$ an input
position and $e_j$ a bit string, with $\sigma_j(r) = K - j$ for all $j \le K$.
Suppose that for every $j < K$, from label $c$ with registers $\sigma_j$ modified
so that register $r$ holds $K - j - 1$, and $p_j$ symbols read, $P$ reaches label
$h$ with registers $\sigma_{j+1}$ and $p_{j+1}$ symbols read within $b$ steps
while appending $e_j$ to its output.  Then from label $h$ with registers
$\sigma_0$ and $p_0$ symbols read, $P$ reaches label $z$ with registers
$\sigma_K$ and $p_K$ symbols read within $K(b+2)+1$ steps while appending the
concatenation $e_0 e_1 \cdots e_{K-1}$.
\end{theorem}

\begin{definition}[Template micro-operation]
\lean{Complexity.CounterProg.MOp}
\label{Complexity.CounterProg.MOp}
\leanok
For $R \in \mathbb{N}$, a template micro-operation over $R$ registers is either
$\mathsf{out}(b)$ for a bit $b$, meaning ``print the bit $b$'', or
$\mathsf{pr}(r)$ for a register index $r < R$, meaning ``print the value of
register $r$ in unary''.
\end{definition}

\begin{definition}[Output of a micro-operation]
\lean{Complexity.CounterProg.MOp.exec}
\label{Complexity.CounterProg.MOp.exec}
\leanok
\uses{Complexity.CounterProg.MOp}
Given register values $\rho : \{0,\dots,R-1\} \to \mathbb{N}$, the string printed
by a template micro-operation is defined by: $\mathsf{out}(b)$ prints the
one-symbol string $b$, and $\mathsf{pr}(r)$ prints $1^{\rho(r)}$, the string of
$\rho(r)$ ones.
\end{definition}

\begin{definition}[Instruction of a template at a position]
\lean{Complexity.CounterProg.tmplInstr}
\label{Complexity.CounterProg.tmplInstr}
\leanok
\uses{Complexity.CounterProg.Instr, Complexity.CounterProg.MOp}
Given a list $\tau = (m_0, \dots, m_{n-1})$ of template micro-operations, labels
$\lambda : \mathbb{N} \to \Lambda$ and a continuation label $\nu$, the
counter-program instruction at index $i$ is: $\mathsf{out}(b, \lambda(i+1))$
(print $b$, go to $\lambda(i+1)$) if $i < n$ and $m_i = \mathsf{out}(b)$;
$\mathsf{pr}(r, \lambda(i+1))$ (print register $r$ in unary, go to
$\lambda(i+1)$) if $i < n$ and $m_i = \mathsf{pr}(r)$; and $\mathsf{goto}\ \nu$
if $i \ge n$.
\end{definition}

\begin{theorem}[Running a straight-line template]
\lean{Complexity.CounterProg.goes_tmpl}
\label{Complexity.CounterProg.goes_tmpl}
\leanok
\uses{Complexity.CounterProg.Goes, Complexity.CounterProg.Goes.mono, Complexity.CounterProg.Goes.trans, Complexity.CounterProg.MOp, Complexity.CounterProg.MOp.exec, Complexity.CounterProg.goes_goto, Complexity.CounterProg.goes_out, Complexity.CounterProg.goes_pr, Complexity.CounterProg.tmplInstr}
\difficulty{5}
Let $P$ be a counter program on a binary input $x$, let
$\tau = (m_0, \dots, m_{n-1})$ be a list of micro-operations, each either
``print a bit $b$'' or ``print register $r$ in unary'', and let
$\lambda : \mathbb{N} \to \Lambda$ and $\nu \in \Lambda$ be labels.  Suppose that
for each $i < n$ the instruction of $P$ at $\lambda(i)$ performs $m_i$ and then
goes to $\lambda(i+1)$, and the instruction at $\lambda(n)$ is
$\mathsf{goto}\ \nu$.  Then for all registers $\rho$ and input positions $p$,
from label $\lambda(0)$ with registers $\rho$ and $p$ symbols read the program
reaches label $\nu$ with registers $\rho$ and $p$ symbols read within $n + 1$
steps, appending $w_0 w_1 \cdots w_{n-1}$ to its output, where $w_i$ is $b$ if
$m_i$ prints the bit $b$ and $w_i = 1^{\rho(r)}$ if $m_i$ prints register $r$.
\end{theorem}

\begin{definition}[Linear expression in the registers]
\lean{Complexity.CounterProg.LinE}
\label{Complexity.CounterProg.LinE}
\leanok
For $R \in \mathbb{N}$, a linear expression over $R$ registers is a pair
$(L, c)$ consisting of a finite list $L$ of register indices in
$\{0, \dots, R-1\}$, repetitions allowed, and a constant $c \in \mathbb{N}$; it
stands for the sum of the listed registers plus $c$.
\end{definition}

\begin{definition}[Value of a linear expression]
\lean{Complexity.CounterProg.LinE.val}
\label{Complexity.CounterProg.LinE.val}
\leanok
\uses{Complexity.CounterProg.LinE}
Given register values $\rho : \{0,\dots,R-1\} \to \mathbb{N}$ and a linear
expression $(L, c)$ with $L = (r_1, \dots, r_k)$, its value is the natural number
$\rho(r_1) + \dots + \rho(r_k) + c$.
\end{definition}

\begin{definition}[Micro-operations printing a linear expression]
\lean{Complexity.CounterProg.linOps}
\label{Complexity.CounterProg.linOps}
\leanok
\uses{Complexity.CounterProg.LinE, Complexity.CounterProg.MOp}
For a linear expression $(L, c)$ over $R$ registers with $L = (r_1, \dots, r_k)$,
this is the list of template micro-operations $\mathsf{pr}(r_1), \dots,
\mathsf{pr}(r_k)$ (print each listed register in unary) followed by $c$ copies of
$\mathsf{out}(1)$ (print the bit $1$).
\end{definition}

\begin{theorem}[A linear expression prints its value in unary]
\lean{Complexity.CounterProg.flatMap_exec_linOps}
\label{Complexity.CounterProg.flatMap_exec_linOps}
\leanok
\uses{Complexity.CounterProg.LinE, Complexity.CounterProg.LinE.val, Complexity.CounterProg.MOp, Complexity.CounterProg.MOp.exec, Complexity.CounterProg.linOps}
\difficulty{4}
Let $\rho : \{0,\dots,R-1\} \to \mathbb{N}$ be register values and $(L, c)$ a
linear expression with $L = (r_1, \dots, r_k)$.  Concatenating the strings
printed under $\rho$ by the micro-operations $\mathsf{pr}(r_1), \dots,
\mathsf{pr}(r_k)$ followed by $c$ copies of $\mathsf{out}(1)$, where
$\mathsf{pr}(r)$ prints $1^{\rho(r)}$ and $\mathsf{out}(1)$ prints $1$, gives
exactly the unary string $1^{v}$ with $v = \rho(r_1) + \dots + \rho(r_k) + c$.
\end{theorem}

\begin{definition}[Micro-operations printing a fixed bit string]
\lean{Complexity.CounterProg.bitsOps}
\label{Complexity.CounterProg.bitsOps}
\leanok
\uses{Complexity.CounterProg.MOp}
For a bit string $w = w_1 w_2 \cdots w_n$, this is the list of template
micro-operations $\mathsf{out}(w_1), \dots, \mathsf{out}(w_n)$, each printing one
bit of $w$.
\end{definition}

\begin{theorem}[A fixed bit string prints itself]
\lean{Complexity.CounterProg.flatMap_exec_bitsOps}
\label{Complexity.CounterProg.flatMap_exec_bitsOps}
\leanok
\uses{Complexity.CounterProg.MOp, Complexity.CounterProg.MOp.exec, Complexity.CounterProg.bitsOps}
\difficulty{3}
For any register values $\rho$ and any bit string $w = w_1 \cdots w_n$,
concatenating the strings printed under $\rho$ by the micro-operations
$\mathsf{out}(w_1), \dots, \mathsf{out}(w_n)$, where $\mathsf{out}(b)$ prints the
single bit $b$, gives back $w$.
\end{theorem}

\begin{theorem}[Machine simulation of a bounded run]
\lean{Complexity.CounterProg.sim_run}
\label{Complexity.CounterProg.sim_run}
\leanok
\uses{Complexity.CounterProg.St, Complexity.CounterProg.enc, Complexity.CounterProg.run, Complexity.CounterProg.run_of_halted, Complexity.CounterProg.run_succ, Complexity.CounterProg.sim_step, Complexity.CounterProg.step, Complexity.CounterProg.step_pos_le, Complexity.CounterProg.step_regs_le, Complexity.CounterProg.toTM, Turing.MultiTapeTM.runFrom, Turing.MultiTapeTM.runFrom_add}
\difficulty{4}
Let $P$ be a counter program over $R$ registers with a finite label set
$\Lambda$, let $\ell_0 \in \Lambda$, let $x$ be a binary input, and let $M$ be the
multi-tape Turing machine compiled from $P$ with start label $\ell_0$, which keeps
each register on its own work tape in unary.  For an abstract state $s$ (label or
halted, registers $\rho_s$, number $\mathrm{pos}(s)$ of input symbols read,
output so far) let $\mathrm{enc}(s)$ be the configuration of $M$ in the control
state for the label of $s$ (or halted), with the input head on the next unread
symbol, work tape $r$ holding $1^{\rho_s(r)}$ with its head on cell
$\rho_s(r)$, and output equal to that of $s$.  Let $B \in \mathbb{N}$.  Then for
every $t \in \mathbb{N}$ and every abstract state $s$ with
$\mathrm{pos}(s) \le |x|$ and $\rho_s(r) + t \le B$ for every register $r$,
there is $t' \le t(2B+3)$ such that running $M$ for $t'$ steps from
$\mathrm{enc}(s)$ yields $\mathrm{enc}(s_t)$, where $s_t$ is the state reached
by $t$ steps of $P$ on $x$ from $s$.
\end{theorem}

\begin{theorem}[A counter program is a Turing machine]
\lean{Complexity.CounterProg.exists_tm}
\label{Complexity.CounterProg.exists_tm}
\leanok
\uses{Complexity.CounterProg.St, Complexity.CounterProg.enc, Complexity.CounterProg.init, Complexity.CounterProg.run, Complexity.CounterProg.sim_run, Complexity.CounterProg.toTM, Turing.Cfg.init, Turing.FinTM.ComputesInTime, Turing.FinTM.computesInTime_iff, Turing.MultiTapeTM.initCfg, Turing.MultiTapeTM.runFrom, Turing.MultiTapeTM.runFrom_add, Turing.MultiTapeTM.runFrom_of_halt, Turing.UnaryTape.ones_zero}
\difficulty{5}
Let $P$ be a counter program over $R$ registers with a finite label set
$\Lambda$ and start label $\ell_0$, let $x$ be a binary input and
$t \in \mathbb{N}$.  Suppose that $t$ steps of $P$ on $x$, started from the
initial state (label $\ell_0$, all registers $0$, nothing read, empty output),
reach a halted state, and let $y$ be its output.  Then the multi-tape Turing
machine compiled from $P$ with start label $\ell_0$, which keeps each register
on its own work tape in unary, computes $y$ on input $x$ in time $t(2t+3)$:
after $t(2t+3)$ steps from its initial configuration on $x$ it is halted with
output $y$.
\end{theorem}

\begin{theorem}[Input position grows at most one per step]
\lean{Complexity.CounterProg.run_pos_le}
\label{Complexity.CounterProg.run_pos_le}
\leanok
\uses{Complexity.CounterProg.St, Complexity.CounterProg.run, Complexity.CounterProg.run_succ, Complexity.CounterProg.run_zero, Complexity.CounterProg.step}
\difficulty{4}
Let $P$ be a counter program on a binary input $x$, let $s$ be an abstract state
of $P$ in which $\mathrm{pos}(s)$ input symbols have been read, and let
$t \in \mathbb{N}$.  Then after $t$ steps of $P$ on $x$ from $s$ the number of
input symbols read is at most $\mathrm{pos}(s) + t$.
\end{theorem}

\begin{theorem}[A step only appends to the output]
\lean{Complexity.CounterProg.step_out}
\label{Complexity.CounterProg.step_out}
\leanok
\uses{Complexity.CounterProg.St, Complexity.CounterProg.step}
\difficulty{3}
Let $P$ be a counter program on a binary input $x$ and $s$ an abstract state of
$P$ with output $o$.  Then there is a bit string $e$ such that the output after
one step of $P$ on $x$ from $s$ is $o\,e$.
\end{theorem}

\begin{theorem}[A run only appends to the output]
\lean{Complexity.CounterProg.run_out}
\label{Complexity.CounterProg.run_out}
\leanok
\uses{Complexity.CounterProg.St, Complexity.CounterProg.run, Complexity.CounterProg.run_succ, Complexity.CounterProg.run_zero, Complexity.CounterProg.step, Complexity.CounterProg.step_out}
\difficulty{4}
Let $P$ be a counter program on a binary input $x$, $s$ an abstract state of $P$
with output $o$, and $t \in \mathbb{N}$.  Then there is a bit string $e$ such
that the output after $t$ steps of $P$ on $x$ from $s$ is $o\,e$.
\end{theorem}

\begin{theorem}[Output growth of one step]
\lean{Complexity.CounterProg.step_out_le}
\label{Complexity.CounterProg.step_out_le}
\leanok
\uses{Complexity.CounterProg.St, Complexity.CounterProg.step}
\difficulty{3}
Let $P$ be a counter program on a binary input $x$, $K \in \mathbb{N}$, and $s$
an abstract state of $P$ with output $o$ in which every register holds at most
$K$.  Then the output after one step of $P$ on $x$ from $s$ has length at most
$|o| + K + 1$.
\end{theorem}

\begin{theorem}[Output length bound from the initial state]
\lean{Complexity.CounterProg.run_init_out_le}
\label{Complexity.CounterProg.run_init_out_le}
\leanok
\uses{Complexity.CounterProg.init, Complexity.CounterProg.run, Complexity.CounterProg.run_add, Complexity.CounterProg.run_regs_le, Complexity.CounterProg.run_zero, Complexity.CounterProg.step, Complexity.CounterProg.step_out_le}
\difficulty{4}
Let $P$ be a counter program on a binary input $x$, $\ell_0$ a label and
$t \in \mathbb{N}$.  Starting from the initial state at $\ell_0$ (all registers
$0$, nothing read, empty output), the output after $t$ steps of $P$ on $x$ has
length at most $t(t+1)$.
\end{theorem}
